% !TeX root = ../queens.tex
\section{Why the local records suffice}\label{sec:sufficiency}

Section~\ref{sec:local} described a calculation on states. We now
show that, when the state comes from the board, one of its branches
carries out the actual greedy step. Two things could go wrong: an
upper queen whose symbol is not stored could affect an attack test or
the row count, or the calculation could request a symbol that the
board has not yet determined. The bounds in this section rule out
both.

The argument concerns a single column $n$: it assumes that the records
before column $n$ satisfy the bounds below. Section~\ref{sec:finite}
checks that the state before column $30$ satisfies them and that every
successor does too, and Section~\ref{sec:identification} combines these
checks with the result of this section by induction on $n$.

\begin{condition}[Bounds on the stored records]\label{cond:state}
The records satisfy
\begin{equation}\label{eq:state-bounds}
 w\le4,\qquad -4\le z\le5,\qquad R,D\subseteq[1\dts4].
\end{equation}
\end{condition}

Each bound has a specific use. The bounds on $w$ and $D$ give the
diagonal discrepancy estimate, as explained after
\eqref{eq:local-rank}. The upper bound on $z$, together with the
bounds on $w$ and~$R$, keeps every request within six places of
index $m$. The lower bound on $z$, together with $U(m-1)\ge12$, then
guarantees that every requested symbol is already determined, and it
keeps upper queens before the input history from affecting the
tests.

These bounds alone do not confine the records to finitely many
states. That only finitely many states are reached is a result of
the verification in Section~\ref{sec:finite}.

% Layout only: keep the statement of the lemma on one page.
\Needspace{15\baselineskip}
\begin{lemma}[The records determine the actual step]\label{lem:local-exactness}
Consider the actual board before column~$n$ and its records
(Definition~\ref{def:local-state}). Suppose that
\begin{enumerate}
\item the records satisfy Condition~\ref{cond:state};
\item $U(m-1)\ge12$ and $m+|Q|\le n$; and
\item $\sigma_1\dots\sigma_{n-1}$ follows the history graph.
\end{enumerate}
Follow the branch of the calculation that answers each request with
the symbol of the queen word at the requested index. We call it the
\emph{actual branch}. Then:
\begin{enumerate}
\renewcommand{\labelenumi}{(\roman{enumi})}
\item every requested index is at most $m+6<n$, and each answer
labels an edge leaving the vertex from which it is read, so the
branch never stops;
\item the branch places the actual queen in column $n$, produces
$\sigma_n$, and gives the actual records before column $n+1$;
\item the row reference advances by at most six places.
\end{enumerate}
\end{lemma}

\begin{proof}
As before, write $\kappa=U(m-1)$. Since $z=n-m-\kappa$ on the board,
the hypotheses give
\begin{equation}\label{eq:causality}
 n-m=\kappa+z\ge12-4=8.
\end{equation}

We first show that every requested index is less than $n$. The preliminary
request, to length $z\le5$, reaches at most offset
$4$. A candidate request reaches offset
$\max\{r,\lfloor(z+r)/2\rfloor\}\le4$, since $r\le w\le4$ and
$z+r\le9$. Only the row search goes further. After a lower choice is
inserted, $\widetilde R\subseteq[0\dts4]$, since $R\subseteq[1\dts4]$
and $r\le w\le4$. Of the rows at offsets $5$ and $6$, at least one
has row bit zero, because successive upper rows differ by at least
two (Lemma~\ref{lem:upper}). Hence an unused row occurs by
offset~$6$, and $\mu\le6$. This proves (iii).

Every request is therefore for an index at most $m+6<n$, and the
symbols already in $Q$ have indices less than $n$ by hypothesis~(2). So every bit that the actual branch reads is
already determined by the queens placed before column $n$. In
particular, if the queen placed in column $n$ is upper, its row is
greater than $n$, so it cannot affect a row search that ends by
$m+6$.

Next we show that the branch never stops. Every requested symbol lies in $\sigma_1\dots\sigma_{n-1}$, which
follows the history graph by hypothesis~(3), and the vertex from which
it is read consists of the twelve actual symbols before it. So the
symbol labels an edge leaving that vertex, and the branch never stops.
With the bound on the requests, this proves (i).

It remains to prove (ii). We begin by showing that upper queens outside the
retained columns do not affect the tests. On the board, extend $\Gamma(h)=U(m+h)-\kappa$ to all offsets with
$m+h\ge1$; it agrees with \eqref{eq:C} on the retained offsets. By
\eqref{eq:source-offsets}, an upper queen in column $m+h$ has
relative row $h+\Gamma(h)$ and relative antidiagonal index
$2h+\Gamma(h)$. The tests compare relative rows with the thresholds
$z-1$ and~$z$, both at least $-5$, and relative antidiagonal indices
with $z+r\ge-4$.

An upper queen before the input history has $h\le-13$ and
$\Gamma(h)\le0$, so its relative row is at most $-13$ and its
relative antidiagonal index is at most $-26$. It lies below both
thresholds, so $B(T)$ should count it; $\kappa$ already does, and no
subtraction is needed. Its antidiagonal index is too small to equal
$z+r\ge-4$, so it attacks no candidate.

An upper queen after the stored queue has $h\ge|Q|$ and
$\Gamma(h)\ge1$, since $\Gamma(h)$ counts column $m+h$ itself. When
$\J(z-1)$ and $\J(z)$ are computed, $|Q|\ge z$, so such a queen has
relative row at least $|Q|+1>z$ and is added to neither count. When
candidate $r$ is tested,
$|Q|>\lfloor(z+r)/2\rfloor$, so such a queen has relative
antidiagonal index at least $2|Q|+1>z+r$ and cannot attack the
candidate.

The sets $R,D,A$ detect every attack by a lower queen; the offsets
they discard are irrelevant, as shown in Section~\ref{app:output}. The
row bit $b_{m+r}$, read from $Q$, detects an upper queen in the
candidate's row, and by the previous paragraphs
Proposition~\ref{prop:upper-antidiagonal-test} detects every upper
queen on its antidiagonal. Conversely, all these attacks come from
queens already placed: the retained columns are less than $n$, and an
upper queen in row $m+r<n$ lies in a column less than $m+r$. The
branch therefore chooses the first unattacked lower candidate, as the
greedy rule does. If there is none, the actual queen is upper, and so
is the branch's choice. The two previous paragraphs verify both
hypotheses of Proposition~\ref{prop:upper-row-count} at the thresholds
$n-1$ and $n$, so it gives the actual row bit $b_n$, and hence the
actual symbol $\sigma_n$ in \eqref{eq:output}.

With the actual symbols, the searches find the actual advances $\mu$
and~$\nu$, \eqref{eq:update} gives the actual new sets and the actual
new values of $w$ and $z$, and the word updates keep exactly the
segments of $\sigma$ at their new positions, all with indices less
than $n+1$. With the previous paragraph, this proves (ii).
\end{proof}

Lemma~\ref{lem:local-exactness} does not assert that the output check
passes, or that the successor satisfies Condition~\ref{cond:state}. The exhaustive check
of Section~\ref{sec:closure-check} establishes both.
